% ==============================================================================
% ПАРТИЯ B03-032: DISS-005-B1800 .. DISS-005-B1858
% Глава 3: Унификация имен сигналов и фильтрация нерелевантных
% ==============================================================================

% DISS-005-B1800

% DISS-005-B1801
\subsubsection{Унификация имен аналоговых и дискретных сигналов}\label{sec:ch3_signal_names_unification}

% DISS-005-B1802
Наиболее сложной и многогранной задачей унификации является приведение имен аналоговых и дискретных сигналов к единому стандарту. Исходные осциллограммы содержат имена, присвоенные производителем устройства или инженером при наладке, и это приводит к огромному разнообразию даже для идентичных по физическому смыслу измерений. Например, фазное напряжение А первой секции шин может именоваться «Ua\_S1», «U\_A\_sect1», «Напр. А СШ1» и~т.д. Автоматическая обработка такого разнообразия невозможна без предварительной стандартизации.

% DISS-005-B1803
Процесс унификации имен включает следующие этапы:

% DISS-005-B1804
\paragraph{Анализ существующих имён}
\mbox{}\par\vspace{0.2cm}

% DISS-005-B1805
{\emergencystretch=2em С помощью специальных функций («\allowbreak find\_\allowbreak all\_\allowbreak name\_\allowbreak analog\_\allowbreak signals», «\allowbreak find\_\allowbreak all\_\allowbreak name\_\allowbreak digital\_\allowbreak signals») осуществляется проход по всем «.cfg» файлам в датасете для сбора всех уникальных имен аналоговых и дискретных сигналов и подсчета частоты их встречаемости. На данном этапе было выявлено порядка 3~170 уникальных имен аналоговых сигналов и более 13~514 уникальных имен дискретных сигналов (по состоянию на [14 июня 2025 \commentnote{DISS-005-C0121}{Алексей Евдаков [2]}{Может поменяться}{г.}]; это число может изменяться по мере пополнения датасета).\par}

% DISS-005-B1806
\paragraph{Разработка стандартной структуры именования}
\mbox{}\par\vspace{0.2cm}

% DISS-005-B1807
Для обеспечения однозначности и возможности автоматизированной обработки была разработана стандартизированная структура именования сигналов. Эта структура основана на типовой схеме двухсекционной подстанции. На рисунке (Рисунок~3.5) показаны основные точки измерения напряжений и токов каждой из секций:

\begin{itemize}
% DISS-005-B1808
  \item \textbf{BB VT} (BusBar Voltage \commentnote{DISS-005-C0122}{Алексей Евдаков [2]}{Может нужно и на русский перевести, хотя имена-то у меня в сигналах на английском, так что возможно, и не стоит}{Transformer}): Трансформаторы напряжения на секциях шин, измеряющие фазные напряжения ($U_a$, $U_b$, $U_c$), междуфазные напряжения ($U_{ab}$, $U_{bc}$, $U_{ca}$) и напряжение смещения нейтрали ($3U_0$).
% DISS-005-B1809
  \item \textbf{CT} (Current Transformer): Трансформаторы тока, измеряющие фазные токи ($I_a$, $I_b$, $I_c$).
% DISS-005-B1810
  \item \textbf{CL VT} (Cable Line Voltage Transformer): Трансформаторы напряжения на кабельных линиях, аналогичные по набору измеряемых величин BB VT.
% DISS-005-B1811
  \item \textbf{Residual CT}: Трансформатор тока нулевой последовательности, измеряющий ток $I_n$.
\end{itemize}

\begin{figure}[htbp]
\centering
% DISS-005-B1812
\includegraphics[width=0.7\textwidth]{image90.png}
% DISS-005-B1813
\caption{Типичная схема расположения секций подстанции и точек измерения сигнала.}\label{fig:ch03_substation_measurement_points}
\end{figure}

% DISS-005-B1814
Для аналоговых сигналов принята трехсегментная структура имени, разделенная символом «|», как показано на рисунке (Рисунок~3.6), здесь:

\begin{itemize}
% DISS-005-B1815
  \item \textbf{Имя сигнала, описывающее его тип:} Например, «U» для напряжения, «I» для тока. Для сигналов от нетрадиционных датчиков (катушки Роговского, емкостные делители), требующих особой обработки, к имени типа добавляется суффикс «\_raw» (например, «U\_raw», «I\_raw») для их явного различения.
% DISS-005-B1816
  \item \textbf{Местоположение на схеме и номер секции:} Указывает, где физически расположен датчик. Например, «BusBar-1» (первая секция шин), «CableLine-2» (вторая кабельная линия).
% DISS-005-B1817
  \item \textbf{Фаза сигнала:} «phase: A», «phase: B», «phase: C» для фазных величин, «phase: AB», «phase: BC», «phase: CA» для междуфазных, «phase: N» для величин нулевой последовательности.
\end{itemize}

\begin{figure}[htbp]
\centering
% DISS-005-B1818
\includegraphics[width=0.6\textwidth]{image91.png}
% DISS-005-B1819
\caption{Пример названия аналогового сигнала.}\label{fig:ch03_analog_signal_naming_example}
\end{figure}

% DISS-005-B1820
На рисунке (Рисунок~3.7) более подробно представлена полная структуру возможных стандартизированных имен аналоговых сигналов для одной секции, включая классические и "сырые" (нетрадиционные) измерения напряжений и токов на шинах и кабельных линиях.

\begin{figure}[htbp]
\centering
% DISS-005-B1821
\includegraphics[width=0.95\textwidth]{image92.png}
% DISS-005-B1822
\caption{Структура аналоговых сигналов.}\label{fig:ch03_analog_signals_structure}
\end{figure}

% DISS-005-B1823
Для дискретных сигналов структура именования также стремится к унификации, хотя их разнообразие значительно больше. Учитывается типовая схема расположения выключателей на двухсекционной подстанции, как показано на рисунке (Рисунок~3.8):

\begin{itemize}
% DISS-005-B1824
  \item IFB (Incoming Feeder Breaker): Вводной выключатель.
% DISS-005-B1825
  \item BC (Bus Coupler): Секционный выключатель.
% DISS-005-B1826
  \item BC Bypass: Обходной секционный выключатель (если предусмотрен).
\end{itemize}

% DISS-005-B1827
Стандартизация дискретных сигналов сфокусирована на двух основных категориях:

\begin{itemize}
% DISS-005-B1828
  \item Положение выключателя: Например, «IFB~|~Bus-1~|~open» (вводной выключатель первой секции отключен), «BC~|~Bus-1~|~close~|~phase:~A» (секционный выключатель первой секции, состояние полюса А~-- включен).
% DISS-005-B1829
  \item Команды управления выключателем: Например, «IFB~|~Bus-1~|~command~open» (команда на отключение вводного выключателя первой секции).
\end{itemize}

% DISS-005-B1830
Рисунок~3.9 иллюстрирует примеры стандартизированных имен для дискретных сигналов, связанных с вводными и секционными выключателями. В случае, если в одной осциллограмме встречается несколько дискретных сигналов с идентичными стандартизированными именами (например, несколько команд управления), они объединяются по логическому условию «ИЛИ». Важно отметить, что для секционных выключателей (BC) к имени добавляется номер секции, чтобы учесть возможность наличия нескольких секционных выключателей на подстанциях с более чем двумя секциями.

\begin{figure}[htbp]
\centering
% DISS-005-B1831
\includegraphics[width=0.6\textwidth]{image93.png}
% DISS-005-B1832
\caption{Структура аналоговых сигналов.}\label{fig:ch03_discrete_signals_structure_layout}
\end{figure}

% DISS-005-B1833

\begin{figure}[htbp]
\centering
% DISS-005-B1834
\includegraphics[width=0.95\textwidth]{image94.png}
% DISS-005-B1835
\caption{Структура аналоговых сигналов.}\label{fig:ch03_discrete_signals_structure_diagram}
\end{figure}

% DISS-005-B1836
\paragraph{Создание словарей сопоставления (маппинга)}
\mbox{}\par\vspace{0.2cm}

% DISS-005-B1837
На основе ручного анализа наиболее часто встречающихся в исходном датасете имен сигналов и разработанной стандартной структуры были созданы таблицы сопоставления (маппинги). Эти таблицы хранятся в формате CSV-файлов и служат основой для автоматического переименования.

\begin{itemize}
% DISS-005-B1838
  \item dict\_\allowbreak analog\_\allowbreak names\_\allowbreak vXX.json (или аналогичный CSV) содержит сопоставления для аналоговых сигналов.
% DISS-005-B1839
  \item dict\_\allowbreak discrete\_\allowbreak names\_\allowbreak vXX.json (или аналогичный CSV) содержит сопоставления для дискретных сигналов.
\end{itemize}

% DISS-005-B1840
В этих файлах каждой уникальной комбинации «сырого» имени сигнала (или его ключевых частей) ставится в соответствие стандартизированное (универсальное) имя. Процесс формирования этих словарей является итеративным: по мере обнаружения новых, ранее не учтенных вариантов именования, словари дополняются и корректируются. Отметим, что данный этап практически полностью ручной.

% DISS-005-B1841
\paragraph{Автоматическое переименование сигналов}
\mbox{}\par\vspace{0.2cm}

% DISS-005-B1842
Специальные функции («rename\_\allowbreak analog\_\allowbreak signals» и «rename\_\allowbreak digital\_\allowbreak signals») используют созданные заранее словари сопоставления имён с универсальными, заменяя исходные имена сигналов на стандартизированные.

% DISS-005-B1843
\paragraph{Обработка дубликатов имен аналоговых сигналов}
\mbox{}\par\vspace{0.2cm}

% DISS-005-B1844
Нередки случаи, когда в одной осциллограмме несколько разных по физическому смыслу сигналов имеют одинаковые или очень похожие имена, отличающиеся, например, только номером присоединения, который не всегда отражен в самом имени. Особую сложность представляют полные дубликаты полученных унифицированных имен аналоговых сигналов внутри одного «cfg» файла. Даже после этапов унификации и нормализации, в ходе которых исправлялись многие ошибки именования, некоторое количество дубликатов унифицированных имен аналоговых сигналов внутри одного файла .cfg может сохраняться (статистика по таким случаям приведена в Таблице~2). Для решения этой проблемы была разработана отдельная функция («find\_\allowbreak duplicate\_\allowbreak signal\_\allowbreak names\_\allowbreak in\_\allowbreak cfg»), которая осуществляет поиск данных сигналов и при активации соответствующей функции~-- переименовывает сигналы. Данный алгоритм выполняет следующие шаги:

\begin{itemize}
% DISS-005-B1845
  \item Анализирует все аналоговые сигналы в файле.
% DISS-005-B1846
  \item Если обнаруживаются полные дубликаты имен, то первый встреченный сигнал (по порядку в файле) остается с исходным именем.
% DISS-005-B1847
  \item Для последующих дубликатов этого же имени предпринимается попытка переназначить их на «свободную» секцию. Например, если есть два сигнала «U | Bus-1 | phase: A», и в осциллограмме нет сигналов, относящихся ко второй секции (Bus-2), то второй дубликат может быть переименован в «U | Bus-2 | phase: A».
% DISS-005-B1848
  \item Если свободной существующей секции нет, создается новая виртуальная секция (например, если максимальный номер секции был 2, то дубликат может быть отнесен к секции 3: «U | Bus-3 | phase: A»).
% DISS-005-B1849
  \item Также ведётся логирование всех произведенных переименований.
\end{itemize}

% DISS-005-B1850
Такой подход позволяет сохранить все исходные данные, устранив при этом неоднозначность имен. Что позволяет в будущем использовать эти сигналы без существенных проблем, а также исключает случаи, когда сигналы с разными номиналами оказываются с одинаковым именем и ввиду этого не могут быть корректно нормализованы.

% DISS-005-B1851
Отметим, что процесс унификации имен является одним из наиболее сложных и затратным с точки зрения ручного труда. Огромное разнообразие исходных имен, отсутствие строгих правил их формирования на местах, ошибки и опечатки~-- все это создает значительные трудности. Попытка охватить все варианты каким-либо общим шаблоном оказывается малоэффективной. Полностью ручной проход по десяткам тысяч осциллограмм нереален. Поэтому был выбран комбинированный подход: автоматический сбор статистики по именам, ручное формирование правил для наиболее частых случаев и разработка алгоритмов для обработки частных проблем, таких как дубликаты. Несмотря на это, полной унификации достичь сложно, и некоторые сигналы могут остаться нераспознанными или быть классифицированы неточно. Особенно это касается дискретных сигналов, чье разнообразие на порядок выше, чем у аналоговых, и их назначение часто узкоспециализировано и зависит от конкретного проекта и логики РЗА. По этой причине в текущем датасете основной упор сделан на унификацию аналоговых сигналов и лишь базового набора дискретных сигналов (например, положение и команды управления вводными и секционным выключателями).

% DISS-005-B1852
\subsubsection{Удаление нерелевантных или избыточных сигналов/строк}\label{sec:ch3_removing_irrelevant_signals}

% DISS-005-B1853
Также, после унификации, часть сигналов, не вошедших в стандартный набор «универсальных» сигналов, становится избыточной для задач формирования общего датасета. Поэтому производится:

\begin{itemize}
% DISS-005-B1854
  \item Удаление пустых строк из «.cfg» файлов.
% DISS-005-B1855
  \item Отбор и сохранение в итоговых обработанных файлах только тех сигналов, которые были успешно унифицированы и соответствуют стандартному перечню. Это позволяет уменьшить «шум» в данных и сосредоточиться на анализе ключевых измерений.
\end{itemize}

% DISS-005-B1856
Завершив этапы деперсонализации и унификации, мы получаем набор осциллограмм, которые являются анонимными, имеют стандартизированные имена сигналов и сгруппированы по частотным характеристикам, что создает основу для следующего важного шага~-- нормализации данных.

% DISS-005-B1857

% DISS-005-B1858
